%========================================================================================
%   General Introduction (unnumbered)
%========================================================================================
\pagenumbering{arabic}
\setcounter{page}{1}
\chapter*{General Introduction}
\addcontentsline{toc}{chapter}{General Introduction}
\markboth{General Introduction}{}

Road freight remains one of the main ways goods circulate within Morocco, and the companies that operate it generate a continuous flow of information: which vehicle left which site, when and with what load, how much fuel was consumed, which advance was paid to the driver, and what price was finally invoiced to the client. For a long time this information has been kept on paper or in spreadsheets, which makes it difficult to consolidate, to audit and to protect. The digitalisation of logistics aims at replacing these scattered records with a single, shared information system. As soon as such a system is put online, however, it concentrates data that is commercially sensitive (client identities, tariffs, margins) and becomes exposed to the same threats as any other web application. Security can therefore not be treated as an afterthought of a logistics project.

This report presents the work carried out during our internship at Youssra Trans-Négoce, a freight transport company based in Casablanca. The company needs to record and follow its transport operations, and our mission was to design and build a web application, named Delivery Tracker, that manages clients, drivers and transport logs. The difficulty of the project is not limited to writing the application itself. A business tool that handles financial and operational data must be developed, tested, deployed and operated in a way that limits the introduction of vulnerabilities and that can be reproduced without relying on manual, error-prone steps.

To address this concern we adopted a DevSecOps approach, in which security checks and automation are integrated into every stage of the software lifecycle instead of being applied at the end. Concretely, the application is a Spring Boot REST API secured with JSON Web Tokens and a React single-page frontend, backed by a PostgreSQL database. Its source code is built and tested by GitHub Actions workflows that also run secret detection, static analysis, dependency and container-image scanning. The infrastructure on which it runs is described as code with Terraform and configured with Ansible, and the application itself is packaged as Docker images. The purpose of this report is to describe these choices, to explain why they were made and to assess honestly what they bring and where they remain incomplete.

The report is organised in nine chapters. Chapter~1 presents the host organisation, the problem to be solved, the objectives, the scope and the planning of the project. Chapter~2 reviews the state of the art on DevSecOps and the tools that support it. Chapter~3 analyses the requirements of the application and of its delivery pipeline. Chapter~4 describes the design of the architecture, of the data model, of the security mechanisms and of the infrastructure. Chapter~5 details the implementation of the backend, the frontend, the containers, the infrastructure and the pipelines. Chapter~6 presents the results obtained and evaluates them. Chapter~7 reports the difficulties we met and the solutions we adopted, and Chapter~8 discusses the limits of the work. Chapter~9 concludes the report and proposes directions for future work.
